\chapter{The octant lexicon}

\textbf{Purpose:} Record what an eight-letter reading of a lit set does and, at greater length, what it cannot do, including two conclusions withdrawn when the object being read turned out to be the wrong one. \textbf{Scope:} \texttt{examples/\allowbreak{}00\_blob\_viz\_tools/\allowbreak{}octant\_lex.py}, \texttt{quotient\_coherence.py}, and SHA-256 compressing one block as the instance.

Eight regions, one letter each, and a word that a computation spells as it runs. The alphabet is general. It reads a lit set on a boundary and knows nothing about what lit it. The same eight letters serve any object that can be placed on a sphere.

Stated first, because a reader who takes one thing away should take this: the alphabet carries eight numbers about a source with \(256\) degrees of freedom, and the result that looks like a success is close to free. Neither of those is a defect in the tool. They are the size of the instrument, counted.

\section{The alphabet}

Choose some regions. A state writes one letter per region, the letter being the share of the lit set sitting inside it. A step of the computation rewrites every letter at once, and the numbers tracked across the steps are the word it spells.

\textbf{The regions are free.} They may sit at any origin, inside the object or outside it, they may overlap each other and anything else, and they may take any shape. The same topology rules govern them as govern the rest of the dimensions, and shape does not enter: a region has to exist, and past that its geometry is a drawing convenience carrying nothing.

Two properties do follow from choices, and they are worth keeping apart from the reading itself. Regions tiling without gap or overlap force the letters to add to one, which costs one degree of freedom. Regions that overlap carry no such constraint and can therefore hold one number each.

\section{The instance this chapter uses}

Eight regions cut by the sign of each coordinate. Every one has a trilateral right-angle corner at the origin, all eight corners meet at that single point, and they tile space with no gap and no overlap. On the boundary the same split cuts eight congruent spherical triangles, each with three right angles and each of area \(\pi/2\), and the measured shares sum to \(1.000000000000000\).

This is the convenient case and not the general one. It was picked because it is easy to state, easy to draw and cheap to compute, and every number below inherits the eight and the tiling from that choice.

Nothing above mentions a hash. What lit the set is the caller's business.

\section{The size of the instrument}

A reading to degree \(L\) carries exactly \((L+1)^2\) real numbers about its source, whatever the source is. Set that against the number of degrees of freedom in the source and the shortfall is the blindness, before any measurement is taken.

\begin{center}
\begin{tabular}{@{}rrrrr@{}}
\toprule
degree & coefficients & rank & blind & least singular value \\
\midrule
4 & 25 & 25 & 231 & \(4.495\) \\
8 & 81 & 81 & 175 & \(4.365\) \\
12 & 169 & 169 & 87 & \(3.818\) \\
15 & 256 & 256 & 0 & \(3.474 \times 10^{-3}\) \\
16 & 289 & 256 & 0 & \(1.524 \times 10^{-1}\) \\
\bottomrule
\end{tabular}
\end{center}

Against a 256-bit lit set, a reading to degree eight recovers \(81\) directions and is blind in \(175\). The map reaches the rank its coefficient count allows at every degree below the source count. The shortfall in coefficients accounts for the blindness on its own. No sample size, no precision and no number of beams changes it.

The eight-letter alphabet here is rank \(8\), seven free numbers once the weight is fixed, blind in \(248\) of \(256\) directions.

Degree fifteen is the floor, where \((L+1)^2\) first reaches \(256\). Reaching the floor is not reaching a usable reading: the least singular value collapses to \(3.5 \times 10^{-3}\) there and recovers to \(1.5 \times 10^{-1}\) at sixteen. One degree of headroom is worth \(44\) times in conditioning. A reading degree picked by counting coefficients alone lands on fifteen and reports success.

\section{The object being read}

Everything from here down is one instance: SHA-256 compressing a single block, with its \(256\) working state bits taken as the lit set after each round. The hash is what the instrument was tried on, and the trying is not the subject.

A frame is the eight working words as the round leaves them, with nothing added back. That distinction decides every number below, and an earlier version of this page had it wrong.

The check counts which object it has instead of trusting the caller. A round only copies six of its eight words, \(b,c,d\) taking \(a,b,c\) and \(f,g,h\) taking \(e,f,g\). On the working state all \(378\) carried words match the words they come from and on the digest none do.

\section{The placement under it}

The \(256\) state bits go onto a golden placement. One index of shift on that placement is a rigid screw: turn \(2.399963\) radians, axial slide \(-0.007812\), pitch \(-0.003255258\). The pitch holds for every amount. One helix carries them all, and where a shifted bit lands is fixed once the amount is known.

\section{The word}

Compressing the empty message, read to degree eight. The delta is the share of the reading that moved since the round before.

\begin{center}
\begin{tabular}{@{}rrrrrrrrrr@{}}
\toprule
round & oct0 & oct1 & oct2 & oct3 & oct4 & oct5 & oct6 & oct7 & delta \\
\midrule
1 & 0.142 & 0.126 & 0.126 & 0.102 & 0.102 & 0.150 & 0.157 & 0.094 & 0.0\% \\
2 & 0.122 & 0.174 & 0.122 & 0.078 & 0.130 & 0.113 & 0.113 & 0.148 & 12.9\% \\
3 & 0.171 & 0.111 & 0.120 & 0.137 & 0.103 & 0.103 & 0.145 & 0.111 & 14.0\% \\
4 & 0.111 & 0.103 & 0.128 & 0.120 & 0.154 & 0.120 & 0.103 & 0.162 & 12.8\% \\
8 & 0.156 & 0.117 & 0.141 & 0.133 & 0.117 & 0.078 & 0.102 & 0.156 & 8.5\% \\
16 & 0.135 & 0.143 & 0.151 & 0.151 & 0.095 & 0.079 & 0.111 & 0.135 & 10.6\% \\
24 & 0.137 & 0.137 & 0.085 & 0.111 & 0.120 & 0.128 & 0.137 & 0.145 & 7.5\% \\
32 & 0.095 & 0.127 & 0.143 & 0.095 & 0.127 & 0.143 & 0.151 & 0.119 & 10.1\% \\
48 & 0.136 & 0.127 & 0.127 & 0.085 & 0.127 & 0.102 & 0.153 & 0.144 & 12.4\% \\
64 & 0.150 & 0.129 & 0.114 & 0.121 & 0.129 & 0.114 & 0.143 & 0.100 & 5.9\% \\
\bottomrule
\end{tabular}
\end{center}

The digest at the end is \texttt{e3b0c442...7852b855}, the published value for the empty message. That check confirms the trace is the hash and not something shaped like one.

\section{The delta against its floor}

The delta averages \(11.9\%\) over the first eight rounds and \(7.1\%\) over the last eight.

Two unrelated states of the same weight differ by \(9.4\%\) in this same alphabet. That figure is the redraw level, and \texttt{redraw\_level} computes it by drawing \(200\) pairs of independent states at the matched weight and averaging the delta between them. Without it the delta is a number with no scale.

\begin{center}
\begin{tabular}{@{}lr@{}}
\toprule
quantity & value \\
\midrule
redraw level & \(9.4\%\) \\
delta, first eight rounds & \(11.9\%\), or \(1.27\) times the floor \\
delta, last eight rounds & \(7.1\%\), or \(0.76\) times the floor \\
first round under the floor & 5 \\
weight range across the run & \(108\) to \(140\) lit bits \\
\bottomrule
\end{tabular}
\end{center}

The early rounds move the reading further than two unrelated states do. A reading whose delta clears its own independence level is not reporting counting noise, because counting noise is the level it cleared.

Two cautions on that, both of which bound it.

The fall is a trend in the means and not a monotone series. The per-round deltas over the first eight rounds run \(12.9\), \(14.0\), \(12.8\), \(8.4\), \(11.9\), \(14.9\), \(8.5\) and \(12.0\). Round \(5\) is the first crossing and not the last time the delta sits high.

Naming the move is not part of the measurement. A move carrying mass from octant to octant instead of mixing it will raise the delta above independence, and a rigid relabeling of the indices is such a move. Reading the early excess as the register shift is a reading of why, and it has not been measured. The measurement says only that the excess is present on the working state and absent on the digest.

\section{How much of each change the alphabet could see}

The delta says the reading moved. It does not say how much of the change the reading was able to move for, and those are different questions. A reading is a linear map, two states give the same letters exactly when their difference lies in its kernel, and the part of a change lying there is invisible however large it is.

The visible fraction answers the second question: the share of a difference lying in the row space of the reading, which for eight indicator arms holds the per-octant totals of the change and discards the rest.

\begin{center}
\begin{tabular}{@{}lrr@{}}
\toprule
 & visible fraction & against an unstructured draw \\
\midrule
first eight rounds & \(0.0397\) & \(1.48\) times \\
last eight rounds & \(0.0138\) & \(0.51\) times \\
mean over all 63 pairs & \(0.0278\) & \(1.04\) times \\
\bottomrule
\end{tabular}
\end{center}

Read the split and never the mean. The early changes sit in the directions the alphabet can see above what chance would put there, the late changes sit in its blind directions below what chance would put there, and the mean is those two canceling. Taken alone the mean reports that the reading catches an arbitrary few percent, and that reading of it is an artifact of averaging.

The early to late ratio is \(2.88\), against \(1.68\) for the delta over the same split, on \(8\) pairs at each end. The direction is clear and the sample is small, and the ratio should be quoted with the \(8\) in the same breath.

\textbf{The two routes share no denominator, and one of them was raised as an objection to the other.} The delta divides counts by a weight moving \(4.3\) bits a round, and that moving denominator was offered as the reason the delta's trend might be bookkeeping instead of arrangement. The visible fraction is a ratio of two norms of one difference vector in raw counts and carries no such term. It shows the trend larger. The denominator was damping the effect and not producing it.

The control arm is what licenses reading any of this. Differences with no structure in them must return the fraction the row space predicts, and for differences holding the weight that figure is \(7\) of \(255\) and not \(8\) of \(256\): such a difference sums to zero, the all-ones direction lies inside the row space since the eight indicators add to it, and the difference lives in the hyperplane orthogonal to all-ones. Measured over \(200\) draws the control returns \(0.0267\) against a predicted \(0.0275\). \texttt{quotient\_coherence.py} refuses to report the measurement at all when the control misses.

\section{The reading that was corrected}

An earlier version of this page reported the alphabet as saturated: the delta sitting at the redraw level from round one, no room left to show mixing take hold, and eight letters called too coarse for the question. That was measured on the wrong object.

Frames came from \texttt{digest\_after}, which adds the initial values back into the working state before returning. Every frame was therefore a finalized digest, and the feed-forward is a mixing step applied after the round, hiding the round's own structure. The count gives it away: on the working state all \(378\) carried words survive and on the digest none do.

The saturation was therefore a property of the digest sequence and not of the alphabet. \texttt{rounds\_of} reads the working state by default now and keeps the digest to compare against. The coarseness statement was withdrawn, and the counting statement in the size of the instrument replaces it, since that one holds whatever object is read.

\section{The distinctness result is close to free}

Sixty-four rounds wrote sixty-four distinct signatures, a coherence of \(100\%\), with no collisions. This is not evidence that the letters carry meaning, and it should not be presented as a finding.

The rank says why in one line. Seven free numbers separate sixty-four arbitrary states nearly always, whether or not the seven carry anything about the source. Distinctness here reports the resolution of a float and not a property of SHA-256, and any count of distinct signatures quoted without the rank beside it is quoting the float.

It is computed because the complement would have been informative. A collision is a null permutation, a move this reading cannot see. None occurred, which rules that out and establishes nothing beyond it.

What would be evidence: a signature that predicts the round it came from.

\section{The churn is not stationary}

An earlier version of this page reported the arrangement moving at a constant rate from the first round to the last, and an earlier version of the check asserted first that the delta grows as the mixing takes hold and then that it holds level. It does neither. It falls, and it crosses the floor on the way down. Both readings were taken on the digest, where the feed-forward had flattened the run before the alphabet saw it.

\section{The twist}

Torsion is the reading that does carry position, and it does not bin anything. The twist between round \(1\) and round \(64\) reads \(4.715210\) radians. A phase recovers an angle to \(1.8 \times 10^{-13}\) radians while a magnitude reads a rotation as nothing at all.

That figure sits under the same rank bound as everything else here. A phase recovers the turn it is given; it does not add coefficients to the reading.

\section{What a finer alphabet would need}

Refining the cells raises the coefficient count, and the counting table above is the bar it has to clear. A subdivision leaving more cells than there are points to write into them has bought dead symbols, and one reaching the rank floor has bought a reading whose conditioning collapses.

\section{Running it}

\begin{center}
\begin{tabular}{@{}l@{}}
\toprule
command \\
\midrule
\texttt{python examples/00\_blob\_viz\_tools/octant\_lex.py --check} \\
\texttt{python examples/00\_blob\_viz\_tools/octant\_lex.py --message "abc"} \\
\texttt{python examples/00\_blob\_viz\_tools/octant\_lex.py --read digest} \\
\texttt{python examples/00\_blob\_viz\_tools/octant\_lex.py --top 12} \\
\bottomrule
\end{tabular}
\end{center}

\texttt{--check} holds the trace to the published digest, holds the eight letters to a sum of one, counts the carried words to confirm which object it is reading, and holds the delta to clearing the floor early and falling across the run. \texttt{--read digest} reproduces the withdrawn reading, and its own check asserts that no carried word survives there.
